\chapter{Multiply-Accumulate Unit (fmac)}
\label{ch:fmac}

\section{Introduction}

\sig{fmac} computes the sum of a stream of products,
$\sum_i \mathtt{op\_a}_i\times\mathtt{op\_b}_i$, and returns one rounded result
per stream. Two \sig{unpack} modules split the operands. The \sig{mac} module
multiplies them, adds each product to a two's complement accumulator and
normalizes the sum. \sig{pack} rounds the sum to nearest, ties to even, as
Figure~\ref{fig:fmac} shows.

The value fed back to the accumulator is rounded to odd: the bits shifted out are
ORed into its least significant bit, so that the final rounding can see that the
value is inexact. The output normalization uses the full adder word.

\begin{figure}[H]
\centering
\begin{tikzpicture}[node distance=7mm]
  \node[ugblock] (u) {\sig{unpack}\\(2)};
  \node[ugblock,right=of u,minimum width=44mm] (a) {\sig{mac}\\multiply, accumulate, normalize};
  \node[ugblock,right=of a] (p) {\sig{pack}\\1 cycle};
  \draw[ugarrow] ($(u.west)-(9mm,0)$) node[above right,font=\footnotesize\ttfamily,xshift=-1mm] {op} -- (u);
  \draw[ugarrow] (u) -- (a); \draw[ugarrow] (a) -- (p);
  \draw[ugarrow] (p) -- ($(p.east)+(9mm,0)$) node[above left,font=\footnotesize\ttfamily,xshift=1mm] {res};
\end{tikzpicture}
\caption{\sig{fmac} Block Diagram}
\label{fig:fmac}
\end{figure}

\section{Features}

\begin{itemize}
\item Accumulation of a stream of products, with one product per clock cycle.
\item A latency of 8 clock cycles from the last product.
\item Rounding to odd inside the accumulator, and to nearest, ties to even, at the
  output.
\item Accumulator width set by \sig{LGUARD\_W} and \sig{RGUARD\_W}.
\item Subnormal operands and results.
\item NaN and infinity handling, selected by \sig{SPECIAL\_CASES}.
\end{itemize}

\section{Parameters}

\sig{fmac} has the parameters of Table~\ref{tab:common_params} and those of
Table~\ref{tab:fmac_params}. The accumulator word has
$\mathtt{LGUARD\_W}+2\cdot\mathtt{MAN\_W}+\mathtt{RGUARD\_W}+5$ bits.

\begin{paramtable}{\sig{fmac} Additional Parameters}{tab:fmac_params}
LGUARD\_W & 0 & Guard bits added above the product in the accumulator. \\ \hline
RGUARD\_W & 0 & Guard bits added below the product in the accumulator. \\
\end{paramtable}

\section{Interface Signals}

\sig{fmac} has the signals of Table~\ref{tab:common_sigs}. \sig{valid\_i} is high for
as many clock cycles as the stream has products.

\Needspace*{23\baselineskip}
\section{Instantiation}

The example below instantiates \sig{fmac} for binary32.
Compile it with the files of Table~\ref{tab:files}.

\begin{vcode}
fmac #(
   .DATA_W       (32),
   .EXP_W        (8),
   .SPECIAL_CASES(1),
   .LGUARD_W     (0),
   .RGUARD_W     (0)
) fmac0 (
   .clk    (clk),
   .rst    (rst),
   .valid_i(valid_i),
   .valid_o(valid_o),
   .op_a   (op_a),
   .op_b   (op_b),
   .res    (res)
);
\end{vcode}

\section{Operation}

A \emph{stream} is a run of consecutive clock cycles with \sig{valid\_i} high:

\begin{enumerate}
\item Raise \sig{valid\_i} and apply one pair of operands per clock cycle. The unit
  samples the operands only in these cycles.
\item Lower \sig{valid\_i} after the last pair. This ends the stream.
\item \sig{valid\_o} is high for one clock cycle, 8 cycles after the last cycle with
  \sig{valid\_i} high, whatever the stream length. Read \sig{res} in that cycle.
\end{enumerate}

A stream of one pair returns the rounded product. Keep \sig{valid\_i} low for at
least 2 clock cycles between two streams. The next stream can then start before
the \sig{valid\_o} of the previous one. A single idle cycle is illegal: the result
of the earlier stream is wrong. Figure~\ref{fig:timing_mac} shows a stream of three products.

\begin{figure}[H]
\centering
\begin{tikztimingtable}[timing/.style={x=2.7mm,y=3mm},timing/d/text/.style={font=\scriptsize}]
cycle & 2D{0} 2D{1} 2D{2} 2D{3} 2D{4} 2D{5} 2D{6} 2D{7} 2D{8} 2D{9} 2D{10} 2D{11} \\
\sig{clk} & 12{HL} \\
\sig{valid\_i} & 6H 18L \\
\sig{op\_a}, \sig{op\_b} & 2D{p0} 2D{p1} 2D{p2} 18U \\
\sig{valid\_o} & 20L 2H 2L \\
\sig{res} & 20U 2D{res} 2U \\
\end{tikztimingtable}
\caption{Timing Diagram, Stream of Three Products}
\label{fig:timing_mac}
\end{figure}

Table~\ref{tab:fmac_special} lists the results for special operands when
\sig{SPECIAL\_CASES} is 1. The conditions apply over all the products of a stream.
An exact cancellation gives $+0$.

\begin{twotable}{\sig{fmac} Special Results}{tab:fmac_special}{Products in the stream}{Result}{9cm}
Any operand is NaN, or any product is infinity times zero & NaN \\ \hline
Both $+\infty$ and $-\infty$ products & NaN \\ \hline
Infinite products of one sign only & That infinity \\ \hline
Every product has a zero operand and a negative sign & $-0$ \\
\end{twotable}

\section{Accuracy}
\label{sec:fmac_accuracy}

\sig{fmac} passes all 6,133,248 TestFloat level 1 \sig{f32\_mulAdd} vectors,
applied as a stream of two products, $a \cdot b + c \cdot 1.0$. The regular test
run applies the first 200,000 of them (Section~\ref{sec:testfloat}).

A longer stream is not always correctly rounded, because the value fed back to the
accumulator is rounded to odd. The stream test (Section~\ref{sec:streams}) checks
each result against the exact sum, with the error bound given there. In its regular
run of 6,000 streams, with the default \sig{LGUARD\_W} and \sig{RGUARD\_W}, 4,448
results are correctly rounded and the others are within the bound.
